\section*{Chapter \ref{ch:b2:business-cycles} --- Business Cycles and Shocks}

\begin{solution}{exo:b2:business-cycles:1}
% figdata: vol_summary
2.34 per cent (decade to 1913), 6.20 (to 1950), 1.13 (to 2007), 2.10 (to 2020). The decade to 1950 includes the Second World War and its
aftermath: output collapsed in some countries and recovered fast, Germany's volatility reached 27.29 per cent.
\end{solution}

\begin{solution}{exo:b2:business-cycles:2}
% figdata: moments
\omterm{def:b2:business-cycles:procyclical}{Procyclical}: consumption (0.87), investment (0.91), disposable income (0.73). \omterm{def:b2:business-cycles:procyclical}{Countercyclical}: unemployment ($-0.88$). Acyclical:
government purchases ($-0.06$). Investment is the most volatile, 4.66 times output.
\end{solution}

\begin{solution}{exo:b2:business-cycles:3}
% ledger: b2:stock-watson2002; figdata: moderation
Stock and Watson: 2.7 and 1.6 per cent, a fall of $1 - 1.6/2.7 = 41$ per cent. This chapter's data: 2.47 and 1.58, a fall of $1 - 1.58/2.47 = 36$ per cent. The difference comes from the
vintage of the data and the way annual growth is measured.
\end{solution}

\begin{solution}{exo:b2:business-cycles:4}
% ledger: b2:nber-faq
Because it looks at a range of monthly indicators rather than GDP alone, weighs depth, diffusion and duration together,
and gives gross domestic income the same weight as GDP. The 2001 recession had no two consecutive quarters of falling GDP.
\end{solution}

\begin{solution}{exo:b2:business-cycles:5}
% code: exercises/solutions.py iid_half_life
The response is $(1/3)^t$; it equals 0.5 when $t = \ln 2/\ln 3 = 0.63$ periods: less than one period.
\end{solution}

\begin{solution}{exo:b2:business-cycles:6}
% code: exercises/solutions.py shares
Saved $\alpha\beta = 0.32$, consumed 0.68. With log utility the income and substitution effects of a change in the return to capital
cancel, and with full depreciation the future return is proportional to output; so the planner saves the same share whatever technology is.
\end{solution}

\begin{solution}{exo:b2:business-cycles:7}
% code: exercises/solutions.py irf
1.00, 0.83, 0.53, 0.30: from $(0.5^{t+1} - (1/3)^{t+1})/(0.5 - 1/3)$, for example $(0.25 - 0.111)/0.1667 = 0.83$ in period 1.
\end{solution}

\begin{solution}{exo:b2:business-cycles:8}
With full depreciation and log utility, consumption is a constant share of output, so its percentage deviations equal
output's. With partial depreciation, a temporary rise in output raises wealth by less than income, and a household that
smooths consumption saves most of the temporary gain: consumption moves less than output and investment more.
\end{solution}

\begin{solution}{exo:b2:business-cycles:9}
% code: exercises/solutions.py growth_moments
Consumption: 0.79 times as volatile as output, correlation 0.66. Investment: 5.33 times, correlation 0.82. The ranking is the same as with HP cycles
(0.80 and 4.66), but the correlations are lower: growth rates keep quarter-to-quarter noise that the HP cycle averages away.
\end{solution}

\begin{solution}{exo:b2:business-cycles:10}
With $k_{t+1} = \alpha\beta y_t$, $c_t = (1 - \alpha\beta)y_t$. The right side of the \omterm{def:b2:intertemporal-choice:euler-equation}{Euler equation} is $\beta E_t[\alpha y_{t+1}/k_{t+1} \times 1/c_{t+1}] = \beta E_t\bigl[\frac{\alpha y_{t+1}}{\alpha\beta y_t(1 - \alpha\beta)y_{t+1}}\bigr] = \frac{1}{(1 - \alpha\beta)y_t} = 1/c_t$: the
realised $y_{t+1}$ cancels, so the equation holds whatever $a_{t+1}$ turns out to be.
\end{solution}

\begin{solution}{exo:b2:business-cycles:11}
Iterating, $\hat y_t = \sum_{j=0}^{t}\alpha^{t-j}\rho^j = \alpha^t\sum_j(\rho/\alpha)^j = \alpha^t\frac{(\rho/\alpha)^{t+1} - 1}{\rho/\alpha - 1} = \frac{\rho^{t+1} - \alpha^{t+1}}{\rho - \alpha}$. Summing over $t$: $\frac{1}{\rho - \alpha}\bigl(\frac{\rho}{1 - \rho} - \frac{\alpha}{1 - \alpha}\bigr) = \frac{1}{(1 - \rho)(1 - \alpha)}$.
\end{solution}

\begin{solution}{exo:b2:business-cycles:12}
With $\lambda = 0$ the objective is $\sum(y_t - \tau_t)^2$, minimised at zero by $\tau_t = y_t$. As $\lambda \to \infty$, any trend with non-zero second differences has an unbounded penalty, so in
the limit the second differences are zero: $\tau_t$ is linear in $t$, the least-squares line.
\end{solution}

\begin{solution}{pb:b2:business-cycles:1}
% code: exercises/solutions.py BRENLAND, irf, half_life, steady_k, shares, iid_half_life
\textbf{Part I.}
1. $c_t + k_{t+1} = a_tk_t^{0.3}$.
2. $1/c_t = 0.95E_t[0.3a_{t+1}k_{t+1}^{-0.7}/c_{t+1}]$.
3. Buildings and roads last decades; only some inputs, such as seed, are used up within a year.
4. $k^* = (0.3 \times 0.95)^{1/0.7} = 0.285^{1.43} = 0.166$.
5. $y^* = 0.166^{0.3} = 0.584$.
\textbf{Part II.}
6. $s = 0.3 \times 0.95 = 0.285$.
7. 0.715.
8. $\hat y_t = \hat a_t + 0.3\hat y_{t-1}$.
9. By the same percentage as output, since it is a constant share.
10. Likewise.
\textbf{Part III.}
11. 1 per cent.
12. $(0.64 - 0.09)/0.5 = 1.10$ per cent.
13. $(0.512 - 0.027)/0.5 = 0.97$ per cent.
14. In year 1: the harvest is still 0.8 per cent above normal and capital, built from year 0's high output, is 1 per cent above, so output is $0.8 + 0.3 \times 1 = 1.1$.
15. 0.52 per cent.
16. Exactly like output: 1, 1.10, 0.97, \ldots{} per cent.
\textbf{Part IV.}
17. The peak is 1.10; half is 0.55; the response is 0.65 in year 4 and 0.52 in year 5, so five years.
18. $\ln 2/\ln(1/0.3) = 0.58$ years.
19. Almost all of it comes from the persistence of the harvest shock; capital adds a little hump and spreads it.
20. $1/(0.2 \times 0.7) = 7.14$.
21. $1/(1 - 0.8) = 5$ from the shock, times $1/(1 - 0.3) = 1.43$ from capital.
22. It peaks in year 2, at 1.28, and is below half the peak after 17 years.
23. \omterm{def:b2:business-cycles:technology-shock}{Technology shocks} almost as persistent as output itself: the model's own propagation adds little.
24. Brenland's output responds to a 1 per cent harvest shock by $(0.8^{t+1} - 0.3^{t+1})/0.5$ per cent: 1, 1.10, 0.97, 0.80, 0.65, 0.52; it falls below half its peak after five years, and almost all
the persistence comes from the shock.
\end{solution}
